% !PRINTED 465
\weberpage{465}

Wir wollen eine Gleichung eine \gesperrt{Normalgleichung nennen, wenn sie irreducibel ist, und die Eigenschaft hat, dass alle ihre Wurzeln rational (in $\Omega$) durch eine von ihnen ausgedrückt werden können}. Dann ist ein primitives Element eines Normalkörpers $\mu^{\mathrm{ten}}$ Grades Wurzel einer Normalgleichung $\mu^{\mathrm{ten}}$ Grades. Aber es gilt auch das Umgekehrte, dass nämlich, wenn $\varrho$ eine Wurzel einer Normalgleichung ist, $\Omega(\varrho)$ ein Normalkörper gleichen Grades ist. Denn ist $\varrho_0$ die Wurzel der Normalgleichung, durch die alle anderen rational ausgedrückt sind, so ist sicher $\Omega(\varrho_0)$ ein Normalkörper. Mit diesem ist aber $\Omega(\varrho)$ conjugirt und also identisch. Es folgt daraus noch, dass bei einer Normalgleichung jede Wurzel nicht nur durch eine bestimmte unter ihnen, \gesperrt{sondern durch jede beliebige rational ausdrückbar ist}.

Man kann nun auf folgendem Wege beliebige algebraische Körper auf Normalkörper zurückführen.

Es sei
\weq{1}{F(x) = 0}
eine beliebige reducible oder irreducible Gleichung in $\Omega$ vom $m^{\mathrm{ten}}$ Grade, von der wir nur voraussetzen wollen, \gesperrt{dass sie keine mehrfachen Wurzeln habe}. Ihre Wurzeln seien
\weq{2}{\alpha,\ \alpha_1 \ldots \alpha_{m-1}.}
Man erhält daraus $m$ algebraische Körper
\weq{3}{\Omega(\alpha),\ \Omega(\alpha_1) \ldots \Omega(\alpha_{m-1}),}
und man kann offenbar die Function $F(x)$ so wählen, dass unter den Körpern (3) irgend welche gegebene algebraische Körper über $\Omega$ vorkommen. Wir nennen den aus allen Grössen (2), d.\,h. aus allen Wurzeln der Gleichung (1) abgeleiteten Körper über $\Omega$
\weq{4}{N = \Omega(\alpha, \alpha_1 \ldots \alpha_{m-1})}
den \gesperrt{Galois'schen Körper der Gleichung} $F(x) = 0$.

Ist $F(x)$ irreducibel, so sind die Körper (3) die mit $\Omega(\alpha)$ conjugirten Körper. In diesem Falle soll der Körper $N$ die \gesperrt{Norm eines jeden der Körper} $\Omega(\alpha), \Omega(\alpha_1) \ldots \Omega(\alpha_{m-1})$ heissen.

Ist $\Omega(\alpha)$ ein Normalkörper, so ist er mit seiner Norm identisch. Im allgemeinen Falle ist nachzuweisen, \gesperrt{dass $N$ ein Normalkörper ist}. Wir wählen ein primitives Element $\varrho$ des Körpers $N$ (nach \S.\,143) und können dann den Körper $N$ auch
%CONT>>
